\par\bigskip
\Needspace{8\baselineskip}
\SourceLeaf{424}
\begin{center}
\pdfbookmark[2]{Gauss an Olbers, 13. Dezember 1810}{pallas-brief-18101213}
\textsc{Gauss} an \textsc{Olbers}, 13. Dezember 1810.
\end{center}
\noindent\hspace*{\parindent}Meine beiden letzten Briefe haben sich mit Ihrem gekreuzt. Seitdem habe ich meine Arbeit über die Störungen der Pallas durch Jupiter abgeschlossen und eine Übereinstimmung erzielt, die meine kühnsten Erwartungen übertrifft. Hier sehen Sie die Abweichungen:
\begin{center}\begingroup\small\setlength{\tabcolsep}{4pt}
\begin{tabular}{lrrrrrr}
& 1803 & 1804 & 1805 & 1807 & 1808 & 1809\\
Mittl. Länge & \(+1''{,}3\) & \(-3''{,}8\) & \(+3''{,}9\) & \(-3''{,}3\) & \(+3''{,}2\) & \(-1''{,}4\)\\
Helioz. Breite & \(-1''{,}0\) & \(+4''{,}1\) & \(+6''{,}0\) & \(+3''{,}9\) & \(-16''\) & \(-3''{,}9.\)\\
\end{tabular}\endgroup\end{center}
\par
\noindent\hspace*{\parindent}Bei der Berechnung der Störungen hatte ich die Elemente zugrunde gelegt, die Sie auf den beiliegenden Blättern\HistoricalNoteEight{Band VI, S. 63.]} auf S. [1971] als Nr. IV finden. Ich habe begonnen, die gesamte Rechnung nochmals zu wiederholen. Dazu teile ich den Zeitraum von 1803 bis 1811 in sieben oder acht Perioden und lege jeder davon die besonderen Elemente zugrunde, die sich aus meiner ersten Rechnung ergeben haben. Aber wie werden Sie sich mit mir wundern, wenn ich Ihnen sage, dass meine Ephemeride in der Oktoberausgabe der M. C. jetzt um mehr als \(1^{0}\) falsch ist, sofern ich mich auf die Ergebnisse der Störungsrechnungen verlassen kann? Ich glaube zwar, meiner Sache sicher zu sein, aber die Rechnung ist so umfangreich, dass ich erst ganz beruhigt sein werde, wenn Pallas wiedergefunden ist. Dafür wünsche ich mir jetzt ungeduldig günstiges Wetter. Sollten Sie in Bremen mehr Glück haben, schauen Sie sich dort doch auch einmal nach ihr um. Nach den gestörten Elementen sollte die AR am 19. Dez. um \(1^{0}10'\) und am 10. Jan. um \(1^{0}40'\) kleiner und die südliche Deklination um \(12'\,.\,.\,.\,10'\) größer sein als in der Ephemeride. Bei der {\AstroSymbols ☌} 1809 mit {\AstroSymbols ♃} wurde {\AstroSymbols ⚴} besonders stark gestört. Bestätigt sich meine Rechnung auch hier zu meiner Zufriedenheit, werde ich vielleicht Juno und Vesta auf ähnliche Weise behandeln.
\par
\noindent\hspace*{\parindent}Die Abhandlung, deren Anzeige ich Ihnen hier beilege, wird bald gedruckt werden.\quad.\quad.\quad.\quad.\quad.
\par
